\section{Related Work}
\label{sec:related}

\subsection{CLIP for Video Action Recognition}
CLIP~\cite{radford2021clip} aligns image and text in a shared embedding space and has been widely adapted to video action recognition.
ActionCLIP~\cite{wang2022actionclip} applies CLIP to video with novel data augmentation and prompt ensembles, establishing a parameter-efficient recipe that keeps the backbone frozen.
X-CLIP~\cite{ni2022xclip} introduces attention-based frame aggregation and video-specific prompts, and achieves strong results on general video recognition.
Vita-CLIP~\cite{wasim2023vitaclip} combines spatial-temporal modules with multimodal prompting and text-adapter alignment, and is among the strongest frozen CLIP baselines.
M2-CLIP~\cite{wang2024m2clip} further couples multimodal prompting with a multi-task objective to adapt CLIP to video in a parameter-efficient way, and VLPA-CLIP~\cite{wang2025vlpa} adapts CLIP with video-language prompting for efficient recognition, while FROSTER~\cite{ju2024froster} shows that a frozen CLIP can serve as a strong teacher for open-vocabulary action recognition.
More recent work strengthens the base representation itself, through stronger video-text encoders~\cite{wang2023internvideo,zhm2025siglip2} and better CLIP alignment~\cite{zhang2025alignclip}, complementing adapter-based adaptation.
Despite these advances, the above methods largely inherit two CLIP conventions: they use only the \texttt{[EOT]} token as the class representation, and they aggregate video tokens uniformly across space and time.
As we analyze in Section~\ref{sec:intro}, both conventions leave the three semantic gaps we target unaddressed, and simply strengthening the video side (as in X-CLIP or Vita-CLIP) does not fix the text-side or the cross-modal gap.
At a higher level, these methods can be organized by which stage of the text--video pipeline they strengthen: prompt-engineering methods~\cite{wang2022actionclip,ni2022xclip,wang2025vlpa} enrich the video side with learnable prompts and aggregation, text-adapter methods~\cite{wasim2023vitaclip} align video features to text, and retrieval-oriented methods~\cite{ju2021clip2video,xue2023clipvip} refine the shared space.
None of them strengthens all three stages simultaneously: the text prototype remains a single token, the video aggregation remains uniform, and the alignment remains directional.
This systematic gap is precisely what MATE-CLIP addresses, through parameter-efficient multimodal adaptation and temporal enhancement with three complementary modules.

\subsection{Video Temporal Modeling}
Video temporal modeling has evolved from 3D convolutions~\cite{tran2015c3d,tran2018i3d} and efficient shift-based designs~\cite{lin2019tsm} to two-pathway architectures~\cite{feichtenhofer2019slowfast} and attention-based transformers~\cite{bertasius2021timesformer,liu2022video_swin,fan2021multiscale,liu2024video_mamba}.
These architectures are trained from scratch on large video corpora, which is costly and still underperforms image-pretrained models on small action datasets.
A recent line of work instead adapts frozen image encoders to video with lightweight modules, such as ST-Adapter~\cite{pan2022st_adapter}, ViT-Adapter~\cite{chen2023vitadapter}, and AIM~\cite{yang2023aim}, drawing on the parameter-efficient transfer recipe of LoRA~\cite{hu2022lora}.
Parameter-efficient designs for CLIP video recognition~\cite{wang2024m2clip,wang2025vlpa} and stronger video-text encoders~\cite{wang2023internvideo,zhm2025siglip2} continue this trend, yet most still model the full token stream and rarely exploit explicit motion correspondence across frames.
However, most adapters still model the full token stream uniformly, so computation is spent on static background as well as on action-critical regions.
Sparse token-selection methods~\cite{li2024video,cao2024rap} reduce this redundancy by focusing on informative tokens.
In particular, RAP~\cite{cao2024rap} selects sparse discriminative patches with a correlated adapter and serves as a strong sparse baseline, but its selection is driven by text-video correlation and does not explicitly exploit motion across frames.
Our MGTA module instead combines \emph{discriminative region selection guided by motion} with \emph{implicit motion-correspondence attention}. In the spirit of deformable attention~\cite{dai2017deformable}, keys and values are sampled from query-dependent offset positions; crucially, however, MGTA predicts a per-token 3D offset that acts on the true spatial--temporal grid, so the offset encodes a \emph{local} correspondence---``this part moves from here to there''---rather than a global uniform shift, which distinguishes it from shift-based temporal models~\cite{lin2019tsm}. The sparse branch thereby \emph{tracks how action-critical regions move over time} and reads the features at their predicted new locations, acting as a lightweight, learnable motion model on top of the frozen encoder.
Unlike token-pruning methods that remove tokens permanently, MGTA performs per-frame selection, so a moving part is re-located in each frame rather than dropped after its first appearance. And unlike dense 3D convolutions~\cite{tran2018i3d,feichtenhofer2019slowfast}, which apply the same kernel everywhere, MGTA spends its capacity only on the regions that actually change, which is both more accurate for subtle motion and cheaper at inference. Compared with the sparse baseline RAP~\cite{cao2024rap}, whose selection is driven by text--video correlation rather than by the dynamics of the video itself, MGTA makes the selection \emph{motion-aware} and couples it with an explicit motion-correspondence branch, which is what allows it to capture how an action changes rather than only where it is.

\subsection{Cross-Modal Alignment}
Aligning vision and language features is central to CLIP and its successors~\cite{li2021blip,li2023blip2}, which bootstrap image-text alignment with auxiliary objectives and frozen encoders.
In video recognition, text-guided alignment~\cite{wasim2023vitaclip,ni2022xclip} aligns video features to text prototypes, and video-text retrieval methods~\cite{ju2021clip2video,xue2023clipvip} refine the shared space with temporal and cross-modal constraints.
More recent video-text encoders~\cite{wang2023internvideo,zhm2025siglip2} and CLIP alignment objectives~\cite{zhang2025alignclip} strengthen this alignment at the representation level, yet in closed-set classification the text directions remain frozen and are rarely corrected toward the actual video class centers.
A key challenge is that direct alignment can disturb already-converged features. Gradient-based alignment toward text directions tends to pull video features away from their natural cluster structure, and the reverse happens when video prototypes are pushed into the text directions.
Prior alignment methods typically train a projection from video to text or refine the text encoder itself. Both choices modify a whole stream of representations and therefore risk un-learning the discriminative structure that the frozen space already encodes.
Our BFEM module instead enhances each modality with the other, using bounded, zero-initialized gates (identity start) and \emph{soft feature enhancement} in both directions: the video enhances each text direction with a soft weight, and the text prototype enhances the video feature, both as convex combinations, while keeping the frozen CLIP backbone intact.
Because the gates open gradually during training, the enhancement starts from the already-converged baseline and only refines the features, rather than re-training the frozen features from scratch.
Unlike projection-based alignment or hard-additive reconstruction, every enhanced feature stays on the segment between the original feature and its modal counterpart, so the enhancement can never drag a representation outside the shared embedding ball.
This distinguishes BFEM from methods that map the video into the text side without ever refining the text directions.
